\section{Diskussion}

\subsection{Zusammenführung der Befunde}

Die einzelnen Analysen ergeben ein konsistentes Bild, das sich entlang der Achse von der Korrelation zur Kausalität ordnen lässt. Die erste Forschungsfrage kann bejaht werden. Sentimentbezogene Informationen sind bereits in den \glspl{token}-Repräsentationen vor der \gls{transformer}-Verarbeitung enthalten und lassen sich mit einer linearen Probe zu 97 Prozent trennen. Die zweite Forschungsfrage beantwortet der Logit Lens. Positive und negative Eingaben besitzen in den frühen Schichten sehr ähnliche Repräsentationen, und die Trennung baut sich schrittweise über den Residual Stream auf. Sie erreicht ihre höchsten Werte in den oberen Schichten zwischen etwa Schicht 17 und Schicht 23. Die dritte Forschungsfrage beantwortet die \gls{attention}-Analyse. Mehrere \glspl{attentionhead} fokussieren auf sentimenttragende Wörter, am stärksten in Schicht 6 und Kopf 3.

Diese Befunde lassen sich auf die Hypothesen beziehen. Der beschreibende Teil von Hypothese H1 wird gestützt, da eindeutig polare Wörter token-nah verarbeitet werden. Dies zeigt sich in der linearen Trennbarkeit der Embeddings sowie in der frühen bis mittleren Aufmerksamkeit auf das Stimmungswort. Der kausale Teil von Hypothese H1, also die Wiederherstellung der Vorhersage durch einen Eingriff am Stimmungswort, steht noch aus. Hypothese H2 zur Lokalisierbarkeit lässt sich erst mit dem Activation Patching abschließend prüfen. Der Logit Lens deutet mit dem Bereich zwischen Schicht 17 und Schicht 23 auf eine mögliche Eingrenzung hin, der Nachweis einer scharf begrenzten und kausal wirksamen Stelle bleibt jedoch offen.

\subsection{Ein scheinbarer Widerspruch}

Auffällig ist ein scheinbarer Widerspruch zwischen den Methoden. Die stärkste direkte Aufmerksamkeit auf sentimenttragende Wörter tritt in den frühen bis mittleren Schichten auf, während die deutlichste Sentiment-Trennung im Residual Stream erst in den oberen Schichten entsteht. Eine plausible Erklärung lautet, dass frühe Köpfe sentimentrelevante Informationen zunächst aus dem Eingabetext einsammeln und in den Residual Stream einspeisen. Die oberen Schichten verdichten diese Informationen anschließend zu einer stabilen Repräsentation und überführen sie in die konkrete Vorhersage. Diese Deutung verbindet die Befunde aus Logit Lens und \gls{attention}-Analyse zu einer gemeinsamen Hypothese über den Verarbeitungsweg.

\subsection{Verhalten und Prompt-Abhängigkeit}

Die verhaltensbasierte Analyse zeigt, dass das Modell sentimentbezogene Informationen nur dann zuverlässig nutzt, wenn der Prompt eine sentimenttragende Fortsetzung wahrscheinlich macht. Der Suffix \texttt{ I feel very} führt mit einer Trefferquote von 75,2 Prozent zu deutlich besseren Ergebnissen als der Suffix \texttt{ Sentiment:} mit 32,1 Prozent. Komplexe sprachliche Phänomene wie Sarkasmus, Ironie, shifter und modality bleiben in beiden Varianten schwierig. Dies steht im Einklang mit der ursprünglichen Motivation des Challenge-Datensatzes, der gerade solche schwierigen Fälle sammelt \cite{barnes2019sentiment}. Eine reine Next-Token-Generierung ersetzt somit keine stabile Sentiment-Klassifikation. Der durchgehend niedrige Wert bei Sarkasmus und Ironie stützt den beobachtbaren Teil von Hypothese H3, wonach das Modell überwiegend der wörtlichen Oberflächen-Polarität folgt. Ob die gemeinte ironische Polarität an keiner Schicht linear ablesbar ist, müsste eine schichtweise lineare Probe auf Ironie-Beispielen klären und bleibt offen.

\subsection{Grenzen der Untersuchung}

Die Ergebnisse sind in mehrfacher Hinsicht einzuordnen. Der Logit Lens verwendet die Ausgabematrix der letzten Schicht. Mittlere Schichten könnten Informationen anders kodieren, sodass einzelne Inhalte übersehen werden. Die lexikonbasierte Bewertung ist grob, da Sentiment auch durch Wörter ausgedrückt werden kann, die nicht im Lexikon enthalten sind. Der Ein-Token-Filter schränkt die betrachtete Wortmenge zusätzlich ein. Alle Analysen beziehen sich auf ein einzelnes Modell mit rund 405 Millionen Parametern, weshalb sich die Befunde nicht ohne Weiteres auf größere Modelle übertragen lassen. Die wichtigste Einschränkung betrifft die Kausalität. Die ersten vier Werkzeuge belegen Korrelationen, aber der kausale Nachweis durch Activation Patching steht noch aus. Die kausalen Anteile der Hypothesen H1 und H2 sind daher noch nicht bestätigt, und auch die schichtweise Prüfung von Hypothese H3 steht aus.
